\begin{theorem}[Bernoulli’s inequality]
Let $x\in\mathbb R$ with $x\ge -1$ and let $n\in\mathbb Z_{\ge 0}$. Then  
\[
(1+x)^{n}\;\ge\;1+nx .
\]

Equality holds exactly in one of the following cases  

\begin{itemize}
\item $n=0$ (any admissible $x$);
\item $x=0$ (any $n\ge 0$);
\item $x=-1$ and $n=1$.
\end{itemize}
\end{theorem}

\begin{proof}
\textbf{1. Degenerate cases.}  

\begin{itemize}
\item \emph{$n=0$.}  $(1+x)^{0}=1$ and $1+0\cdot x=1$, so equality holds.
\item \emph{$x=-1$.}  Then $(1+x)^{n}=0$ and the right–hand side equals $1-n$.  
Since $n\ge0$, $0\ge 1-n$ is equivalent to $n\ge1$, so the inequality is true for all $n\ge1$.  
Equality $0=1-n$ occurs only when $n=1$.
\item \emph{$x=0$.}  Both sides equal $1$, which will also appear in the general equality analysis.
\end{itemize}
Henceforth we assume $n\ge1$ and $x>-1$ (the case $x=-1$ already settled).

\medskip

\textbf{2. Proof for $x\ge0$ (binomial expansion).}  

The Binomial Theorem gives, for any integer $n\ge0$,
\begin{align}
(1+x)^{n}= \sum_{k=0}^{n}\binom{n}{k}x^{k}.
\end{align}
Separating the first two terms,
\begin{align}
(1+x)^{n}=1+nx+\underbrace{\sum_{k=2}^{n}\binom{n}{k}x^{k}}_{S}.
\tag{2}
\end{align}
When $x\ge0$ each binomial coefficient $\binom{n}{k}>0$ and each power $x^{k}\ge0$, hence $S\ge0$. Consequently
\[
(1+x)^{n}\ge 1+nx .
\]
If equality holds then $S=0$. Because every term in $S$ is non‑negative and the coefficients are positive, we must have $x^{k}=0$ for all $k\ge2$, i.e. $x=0$. This yields the equality case $x=0$ (with any $n$).

\medskip

\textbf{3. Proof for the whole interval $[-1,\infty)$ (induction).}  

\emph{Base step $n=0$.}  Already covered in the degenerate cases.

\emph{Inductive hypothesis.}  Assume for a fixed $n\ge0$
\begin{align}
(1+x)^{n}\ge 1+nx \qquad\text{for every }x\ge-1 .
\tag{H\_n}
\end{align}

\emph{Inductive step.}  Since $x\ge-1$ we have $1+x\ge0$, so multiplying \eqref{H\_n} by the non‑negative factor $1+x$ preserves the inequality:
\begin{align}
(1+x)^{n+1}\ge (1+nx)(1+x).
\tag{3}
\end{align}
Expanding the right–hand side,
\begin{align}
(1+nx)(1+x)=1+(n+1)x+nx^{2}.
\tag{4}
\end{align}
Substituting \eqref{4} into \eqref{3} gives
\begin{align}
(1+x)^{n+1}\ge 1+(n+1)x+nx^{2}.
\tag{5}
\end{align}
Because $n\ge0$ and $x^{2}\ge0$, the term $nx^{2}$ is non‑negative. Dropping this non‑negative term yields
\[
(1+x)^{n+1}\ge 1+(n+1)x ,
\]
which is the desired inequality for $n+1$.  

\emph{Equality analysis in the inductive step.}  In \eqref{5} the extra term $nx^{2}$ vanishes precisely when $n=0$ or $x=0$. If $n\ge1$ and $x\neq0$, the term is positive, making the inequality strict. The case $x=-1$ has been handled separately; equality then occurs only for $n=1$.

Thus, by induction, Bernoulli’s inequality holds for every integer $n\ge0$ and every real $x\ge-1$.

\medskip

\textbf{4. Equality condition (combined).}  

From \eqref{2} we have
\[
(1+x)^{n}=1+nx+S,\qquad
S:=\sum_{k=2}^{n}\binom{n}{k}x^{k}\ge0 .
\]
Equality $(1+x)^{n}=1+nx$ is equivalent to $S=0$. Since each coefficient $\binom{n}{k}>0$, each term $\binom{n}{k}x^{k}$ must be zero. Hence either $x=0$ (forcing all terms with $k\ge2$ to vanish) or the sum is empty, i.e. $n\le1$.  

When $n=0$ the statement reduces to $1=1$ for any admissible $x$.  
When $n=1$ the inequality reads $(1+x)\ge1+x$, which is an equality for every $x$; in particular for the endpoint $x=-1$ we obtain the extra equality case $x=-1$, $n=1$.

Collecting all possibilities, equality holds exactly in the three cases listed at the beginning of the theorem.

\end{proof}